\section{Denominators of the local solution}\label{sec:arithmetic}

Our aim is to turn the rational linear forms in \eqref{eq:value-row}
into integer linear forms without losing their smallness. Put $u=z+1$
and let $U(u)$ be the fundamental matrix of \eqref{eq:base-system},
normalized by $U(0)=I$. Thus the solution with initial column $x$ at $-1$
is $U(u)x$. The Taylor coefficients of $U$ record how the derivatives of
a solution depend on its initial values. Since the polynomial row
$C^{[p]}$ has integer coefficients, denominators for these Taylor
coefficients will suffice for our forms.

We use the distinct-product least common multiple from
\cite[Lemma~2]{F26}; here positive indices suffice. The reason for
distinct products will emerge from solving the differential system by
successive integrations: each integration contributes a denominator
equal to a new, strictly larger Taylor degree.

For integers $d\geq1$ and $k\geq0$, define
\begin{equation}\label{eq:Delta-def}
 \Delta_{d,k}
   =\lcm\left\{\prod_{t\in I}t:
        I\subseteq\{1,\ldots,k\},\ |I|\leq d\right\},
 \qquad \Dt_{d,k}=4^k\Delta_{d,k}.
\end{equation}
Empty products and $\lcm(1,\ldots,0)$ are taken to be $1$.

\begin{lemma}\label{lem:lcm}
One has
\begin{equation}\label{eq:lcm-product}
 \Delta_{d,k}=\prod_{j=1}^{d}\lcm(1,\ldots,\lfloor k/j\rfloor).
\end{equation}
In particular, $\Dt_{d,k}$ divides $\Dt_{d,k+1}$, and for fixed $d$,
\begin{equation}\label{eq:D-rate}
 \log\Dt_{d,k}=(\log4+\Ht_d)k+o(k).
\end{equation}
\end{lemma}

\begin{proof}
For a prime $p$, the maximum valuation of a product of at most $d$
distinct integers in $\{1,\ldots,k\}$ is the sum of the $\min\{d,k\}$ largest
valuations in that set. It equals
\[
 \sum_{e\geq1}\min\{d,\lfloor k/p^e\rfloor\}.
\]
The product on the right of \eqref{eq:lcm-product} has the same valuation:
for each $e$, there are exactly $\min\{d,\lfloor k/p^e\rfloor\}$ indices
$j$ with $p^e\leq\lfloor k/j\rfloor$.
Divisibility follows either from the definition or from this formula.
Finally, the prime number theorem gives
$\log\lcm(1,\ldots,K)=K+o(K)$.
Sum this estimate over the fixed set $1\leq j\leq d$ to obtain
\eqref{eq:D-rate}.
\end{proof}

The dependence on $d$ is important. Allowing $d$ arbitrary factors from
$\{1,\ldots,k\}$ would give the larger multiplier
$\lcm(1,\ldots,k)^d$, with logarithm asymptotic to $dk$.
Distinct factors give $\Ht_d k$ instead. Since $\Ht_d$ grows only
logarithmically with $d$, this keeps the arithmetic cost small when the
number of beta values is later increased.

\begin{lemma}\label{lem:denominators}
For the normalized fundamental matrix $U$ defined above,
\[
 \Dt_{d,k}[u^k]U(u)\in M_{d+1}(\Z)
 \qquad(k\geq0).
\]
\end{lemma}

\begin{proof}
The common diagonal term in the equations is removed by an integrating
factor. Put $\eta(u)=(1-u)^{-1/2}$, set $y_i=\eta(u)v_i$ for $i\geq1$,
and leave $y_0=v_0$ unchanged. Since $\eta(0)=1$, this does not change
initial values.
The transformed equations are
\begin{equation}\label{eq:local-gauge}
 v_0'=0,\qquad
 v_1'=-\frac{\eta(u)}{1-u/2}v_0,\qquad
 v_i'=\frac{v_{i-1}}{1-u}\quad(i\geq2).
\end{equation}
Both $\eta(u)$ and $-\eta(u)/(1-u/2)$ have their degree-$j$ coefficient
in $4^{-j}\Z$. For $\eta$, this follows from
\[
 \eta(u)=\sum_{j\geq0}\binom{2j}{j}\frac{u^j}{4^j};
\]
multiplication by $(1-u/2)^{-1}$ preserves the asserted denominator bound.
The coefficients of $(1-u)^{-1}$ are integers.

Each entry of the normalized solution of \eqref{eq:local-gauge} is a
finite sum of iterated integrals of depth at most $d$, and multiplication
by $\eta$ restores $U$. Expand one contribution to a coefficient of degree
$k$. At each integration, the degree increases strictly. Its integration
denominators are therefore distinct integers
\[
 1\leq t_1<\cdots<t_e\leq k,\qquad e\leq d.
\]
Their product divides $\Delta_{d,k}$.
The total degrees supplied by coefficients of the multiplied series,
including the final factor $\eta$, are at most $k$; their powers of $4$
divide $4^k$. This controls each contribution and hence their sum.
The argument includes the column with $v_0(0)=1$, whose first integration
uses the inhomogeneous term displayed in \eqref{eq:local-gauge}.
\end{proof}

\begin{corollary}\label{cor:integer}
For every $p,k,i$ in the construction,
\[
 \Dt_{d,k-1}\frac{Q_{k,i}^{[p]}(-1)}{(k-1)!}\in\Z.
\]
For $k\geq2M+2$ this gives an integer linear form in
$1,\beta(2),\ldots,\beta(2\lfloor d/2\rfloor)$ whose value is
\begin{equation}\label{eq:integer-form}
 \Lambda_{n,p,k}
       =\Dt_{d,k-1}\frac{S_{n,p}^{(k-1)}(-1)}{(k-1)!}.
\end{equation}
\end{corollary}

\begin{proof}
The row $C^{[p]}(-1+u)$ has integral Taylor coefficients.
The coefficient of $u^{k-1}$ in $C^{[p]}(-1+u)U(u)$ is the row
$Q_k^{[p]}(-1)/(k-1)!$, by \eqref{eq:universal-row}.
Use Lemma~\ref{lem:denominators} and
$\Dt_{d,j}\mid\Dt_{d,k-1}$ for $j\leq k-1$.
Substitution of \eqref{eq:H-values} multiplies even columns by integers
and sets odd contributions to zero.
\end{proof}

